\section{Web reconstruction and the polarized first layer}
\label{sec:web-context-v143}
The pencil inverse theorem uses a single irreducible coefficient
module.  For webs, the intrinsic normal cone instead supplies two
projective components that must be recombined.  This supplement keeps
the full web, operator, polarized, and boundary theory.  The statements
below have their original hypotheses and conclusions.

\subsection{The multiplication failure scheme}

Let \(V\) be a complex vector space of dimension
\(e\in\{3,4\}\), put \(p=e(e-1)/2\), and let
\[
 R\subset\Sym^2V
\]
be an \(e\)-dimensional relation space.  Set
\[
 S_R=\Sym^2V/R,\qquad
 B_R=\C\oplus V\oplus S_R,
\]
with multiplication induced by the quotient
\(\Sym^2V\twoheadrightarrow S_R\) and with
\((V\oplus S_R)S_R=0\).  On
\[
 X_R=\Gr(e,V\oplus S_R)
\]
let \(\mathcal K\) be the tautological subbundle.  For every
\(m\ge2\) the multiplication map defines the same zeroth-Fitting
scheme
\(\widehat D_R\) of \eqref{eq:failure-intro}.
Its reduction is the Schubert divisor
\(\widehat\Delta_R\) on which \(\mathcal K\to V\) drops rank.

Regard \(R\) as a web of quadrics on \(\Pj(V^*)\).  On the
invariant open \(G_e^\circ\) used throughout, the web is
basepoint-free and its Jacobian hypersurface
\begin{equation}\label{eq:Y-intro}
 Y_R=V\!\left(\det(\partial q_i/\partial x_j)\right)
 \subset\Pj(V^*)
\end{equation}
is smooth.  For \(e=4\) it is a quartic K3 surface.


\subsection{Recovering the original web}

The web inverse theorem uses the full nonreduced scheme,
including information at the deepest projection-corank stratum.  The
following conclusion holds without an additional genericity condition
and is proved in
Section~\ref{sec:intrinsic-web-reconstruction}.

\begin{theorem}[Intrinsic web reconstruction]
\label{thm:main-web-reconstruction}
For every pair \(R,R'\in G_4^\circ\),
\[
 \widehat D_R\simeq\widehat D_{R'}\quad\Longrightarrow\quad
 R'=gR\ \text{for some }g\in\PGL(V).
\]
Consequently the abstract failure scheme determines \(B_R\)
throughout the smooth Jacobian locus and separates every finite
web-to-K3 packet contained in this locus.
\end{theorem}

The proof first recovers \(\Gr(4,S_R)\) by iterated reduced
singular loci.  Its nonreduced normal cone identifies the two
components of
\[
 \bigwedge^4\Sym^2V=
 \mathbb S_{(4,3,1,0)}V\oplus\mathbb S_{(5,1,1,1)}V.
\]
The second component is the Jacobian covariant; the first is not
retained by the quartic alone.  The two projective components define a
line in the Pluecker space.  The contraction-kernel theorem
(Lemma~\ref{lem:schur-contraction-kernel}) computes the exact exterior
support of the Jacobian component.  For a smooth quartic that support
has dimension ten.  A secant or tangent four-vector on the Grassmannian
has support at most eight.  This excludes every second point and every
infinitesimal ambiguity, so the line meets the Grassmannian in one
reduced point.  In fact three essential variables already suffice.  The global tensor rulings of the normal cone
exclude the transposition ambiguity of an isolated square-matrix
cone.  No extension of an abstract scheme isomorphism to the ambient
Grassmannian is assumed.


The exceptional fibres of the two-component invariant are described in
Theorem~\ref{thm:exceptional-component-fibres}.  A proper intersection
of the component line with the Grassmannian has length at most two.
Its second point can coincide with the first, lie at a discarded pure
endpoint, or give a distinct candidate exchanged by a regular residual
involution.  A positive-dimensional fibre is exactly a Grassmann line
specified by a three-plane inside a five-plane.  The tangent condition
\(Qw_R\in\bigwedge^3R\wedge W\) distinguishes the double-point
case from a reduced secant pair.  All these alternatives are excluded on \(G_4^\circ\) by
Theorem~\ref{thm:global-schur-recombination}.  The ambient obstruction
is confined to the binary-Jacobian boundary, defined invariantly by the
first catalecticant.  Thus the smooth family has no residual secant
involution, flag-line ambiguity, or ramification divisor to resolve.
Outside that family the component-fibre classification is not promoted
to a classification of abstract failure-scheme isomorphisms.


\subsection{The mechanism in arbitrary dimension}

There is a second inverse theorem, for a family that is not confined
to the two-summand plethysm in dimension four.  If \(\dim V=n\),
polarization and the Jacobian determinant give natural maps
\[
 j_n:\det V\otimes\Sym^nV\longrightarrow\bigwedge^n\Sym^2V,
 \qquad C_n:\bigwedge^n\Sym^2V\longrightarrow\det V\otimes\Sym^nV,
 \qquad C_nj_n=2\,\mathrm{id}.
\]
Put \(Q_n=\frac12j_nC_n\) and \(P_n=1-Q_n\).
Theorems~\ref{thm:uniform-contraction}--\ref{thm:uniform-recombination}
prove the exact all-ranks contraction kernel, the resulting support
formula, and the following consequence.  For every \(n\ge4\), on
the nonempty open set of relation spaces whose Jacobians have at
least three essential variables, the pair
\(([P_nw_R],[Q_nw_R])\) determines \(R\).  The corresponding
pencil intersection is a single section even after nonreduced base
change.  The complement \(\ker C_n\) need not be irreducible.

The numerical mechanism is uniform: for \(d<n\) essential variables,
the exterior support is
\(\binom{n+1}{2}-\binom{n-d+1}{2}\), and for \(d=n\) it is full
except for a one-dimensional annihilator of a nondegenerate quadratic
power when \(n\) is even.  The resulting gap above \(2n\) excludes
secants and tangents.  The skew-flattening support criterion itself
is classical \cite{LandsbergOttaviani,Sheridan}; the calculation here
is its restriction to the specified Jacobian component, including
all singular and nondegenerate contraction kernels.  In dimension
four the intrinsic failure scheme supplies the required component
pair, thereby turning the uniform inverse mechanism into
Theorem~\ref{thm:main-web-reconstruction}.  The smooth application
has four essential variables; the algebraic exclusion is stronger
and already applies with three.

\subsection{Intrinsic reconstruction}

On the projection-corank-at-most-one locus let \(D_R\) be the
restriction of \(\widehat D_R\), let \(\NN_R\) be its
nilradical, and let \(E_R\) be the annihilator scheme of the first
nilpotent layer.  The first-layer reconstruction is the following result
proved in Section~\ref{sec:polarization}.

\begin{theorem}[Polarized reconstruction]
\label{thm:main-reconstruction}
For \(R\in G_e^\circ\), put
\[
 \LL_R=\omega_{E_R}\otimes
 \NN_R^{\otimes-(p+e-1)}.
\]
Then
\[
 \bigoplus_{n\ge0}H^0(E_R,\LL_R^{\otimes n})
 \simeq
 \bigoplus_{n\ge0}H^0(Y_R,\OO_{Y_R}(en))
\]
as graded algebras.  In particular, for \(e=4\) the abstract
failure scheme determines the polarized quartic K3
\((Y_R,\OO_{Y_R}(1))\).
\end{theorem}


